\documentclass[10pt,a4paper]{amsart}
\usepackage[margin=19mm]{geometry}
\usepackage{amsmath,amssymb,mathtools}
\usepackage[scr=rsfs,cal=euler]{mathalfa}
\usepackage{xfrac}
\usepackage[unicode,hidelinks,bookmarks=false]{hyperref}
\newcommand{\ie}{i.e.\ }
\newcommand{\cf}{cf.\ }
\newcommand{\resp}{respectively\ }
\newcommand{\Subsection}{\subsection}
\providecommand{\nobreakdash}{\nobreak\hbox{-}}
\setlength{\emergencystretch}{2em}
\multlinegap0pt
\usepackage{amssymb}
\usepackage{enumerate}
\usepackage{xcolor}
\usepackage[all]{xy}
\newtheorem{definition}{Definition}
\newtheorem{remark}{Remark}
\newtheorem{proposition}{Proposition}
\newtheorem{theorem}{Theorem}
\title{Hilbert properties of varieties}
\author{Arno Fehm, Ariyan Javanpeykar}
\date{Source: arXiv:2511.18431v2 (CC BY 4.0)}
\begin{document}
\maketitle

\noindent\emph{Source and licence.} Hilbert properties of varieties. Version arXiv:2511.18431v2 (CC BY 4.0), distributed by arXiv under the Creative Commons Attribution 4.0 International licence, http://creativecommons.org/licenses/by/4.0/, as recorded in the arXiv licence metadata for this exact version and shown on the abstract page https://arxiv.org/abs/2511.18431v2. Full licence notice: \texttt{licenses/arxiv-2511.18431-CC-BY-4.0.txt}.\par\smallskip

\noindent\textit{Editorial scope.} Counted unit: the article's proposition prop:uniruled (source0.tex lines 1424--1432). The article's complete original proof of it is delivered with it (source0.tex lines 1434--1442, one \texttt{proof} environment of 9 lines). No mathematics is written, repaired or re-derived: the statement and the proof are the article's own lines, in the article's own notation, and no retained line is substituted or re-flowed.  Retained uncounted as the source-local context the proof needs: lines 340--347; lines 366--372; lines 545--548; lines 867--881; lines 1008--1019; lines 1092--1102.  Nothing was omitted from any retained span, so the package carries no omission record.  No retained line contains a citation, so the wrapper carries no bibliography and no bibliography entry.  No retained line contains a figure environment, an image-inclusion command or drawing code, so no image is delivered and the package declares no asset dependency.  Editorial formatting only, in addition: an AMS \texttt{amsart} preamble that restates the article's own declarations and macros used by the retained lines, namely the environments \texttt{definition}, \texttt{remark}, \texttt{proposition}, \texttt{theorem} on separate counters, together with the standard packages those lines need.  The retained environments print in this document as Definition 1, Remark 1, Proposition 1, Theorem 1 in order of appearance, whereas the article prints the same items with the numbering of its own full document.  The complete native source of the article as distributed in the arXiv e-print for this exact version is bundled unchanged as \texttt{sources/189652/source0.tex}.
\begin{definition}\label{def:thin}\label{def:HP}
A subset $\Sigma\subseteq X(K)$ is \emph{thin (in $X$)} if there are $n\in\mathbb{N}$
and for each $i\in\{1,\dots,n\}$ a $K$-variety $Y_i$ and a cover $\pi_i\colon Y_i\to X$
with ${\rm deg}(\pi_i)>1$ such that
$\Sigma \setminus \bigcup_{i=1}^n \pi_i(Y_i(K))$
is not Zariski-dense in $X$. 
The variety $X$ has the {\em Hilbert property (``$X$ has HP'')} if $X(K)$ is not thin in $X$.
\end{definition}

\begin{remark}\label{rem:HP}\label{rem:HP_birat}
By definition, HP is a birational property,
so if $X$ is birationally equivalent to a $K$-variety $X'$ (written $X\sim X'$),
then $X$ has HP if and only if $X'$ does.
In Definition \ref{def:HP} one can replace the finite surjective generically \'etale morphisms $\pi\colon Y_i\rightarrow X$ by generically finite generically \'etale dominant rational maps $\pi_i\colon Y_i\dashrightarrow X$ without changing the class of thin sets.
Similarly, one could demand that the $Y_i$ are geometrically integral, since otherwise $\pi_i(Y_i(K))$ is not Zariski-dense in $X$.
\end{remark}

\begin{proposition}\label{prop:WHP_birat}
If $X$ and $X'$ are smooth proper geometrically connected $K$-varieties with $X\sim X'$,
then $X$ has WHP if and only if $X'$ does.
\end{proposition}

\begin{theorem}\label{thm:down} 
Let $f\colon X\rightarrow Y$ be a morphism of normal $K$-varieties.
\begin{enumerate}
\item If $X$ has HP
and $f$ is dominant
and has geometrically irreducible generic fiber,
then  $Y$ has HP.
\item  If $X$ has WHP and
$f$ is smooth, surjective, and has geometrically irreducible generic fiber,
then  $Y$ has WHP.  
\item If $X$ has WHP 
and $f$ is smooth and proper,
then $Y$ has WHP.
\end{enumerate}
\end{theorem}

\begin{theorem}\label{thm:mixed_fibration}\label{thm:fibration}
Let $f\colon X\rightarrow S$ be a dominant morphism of normal $K$-varieties
and $\Gamma\subseteq X(K)$.
Let $\Sigma\subseteq S(K)$ be such that for every $s\in\Sigma$, the fiber $X_s=f^{-1}(s)$
is  integral and  $\Gamma\cap X_s$ is not thin in $X_s$.
\begin{enumerate}
    \item If $\Sigma$ is not thin in $S$, then $\Gamma$ is not thin in $X$.
    \item If $\Sigma$ is not strongly thin in $S$, then $\Gamma$ is not strongly thin in $X$.
\end{enumerate}
In particular, if there is a dense open subset $U\subseteq S$ such that for every $s\in U(K)$ the fiber $X_s$ has HP,
then HP or WHP for $S$ imply HP respectively WHP for $X$.
\end{theorem}

\begin{theorem}\label{thm:curves}
Let $K$ be finitely generated
and let $X$ be a 
 curve over $K$.
 \begin{enumerate}
\item $X$ has HP if and only if $X$ is rational. 
\item  $X$ has WHP if and only if $X$ is rational, or $g_X\leq 1$ and  $X(K)$ is Zariski-dense in $X$.  
\item $X$ has potential HP if and only if $g_X=0$.
\item $X$ has potential WHP if and only if   $g_X\leq 1$.
\end{enumerate}
\end{theorem}

\begin{proposition}\label{prop:uniruled}
    Let $K$ be finitely generated and let $X$ be a surface that is birationally ruled over $K$.
    \begin{enumerate}
        \item $X$ has HP if and only if $q(X)=0$ and $X(K)$ is Zariski-dense in $X$.
        \item $X$ has WHP if and only if $q(X)\leq 1$ and $X(K)$ is Zariski-dense in $X$.
        \item $X$ has potential HP if and only if $q(X)=0$.
        \item $X$ has potential WHP if and only if $q(X)\leq 1$.
    \end{enumerate}
\end{proposition}

\begin{proof}
Since HP and WHP are birational invariants (of smooth projective varieties, see Remark \ref{rem:HP_birat} and Proposition \ref{prop:WHP_birat}),
we can assume that $X=\mathbb{P}^1_K\times C$ for a curve $C$ over $K$ of genus $g_C=q(X)$.
By Theorem \ref{thm:mixed_fibration} and Theorem \ref{thm:down}
applied to the projection $X\rightarrow C$,
$X$ has HP or WHP if and only if $C$ has.
Moreover, $X(K)$ is Zariski-dense in $X$ if and only if $C(K)$ is Zariski-dense in $C$.
Therefore, all claims follow from Theorem \ref{thm:curves} for the curve $C$.
\end{proof}

\end{document}
